% Solution: Gibbons (1992) Problem 1.14. Input by \solitem.
% New in the Lecture 3 practice list. Check: check_L03.py.
\begin{exsol}{Gibbons Problem 1.14}{51}{Advanced}{4}
\asked Show that Proposition~B of Appendix~1.1.C, which says that a pure-strategy Nash equilibrium
survives iterated elimination of strictly dominated strategies, holds for mixed equilibria too:
every strategy played with positive probability in a mixed-strategy Nash equilibrium survives IESDS.

\method Induction on the rounds of elimination, with the principle of indifference (Lecture~3,
slides~21--26) as the key step; the pure-strategy case is the same argument (Lecture~2, slide~16).
Notation: $S_i^0 = S_i$ and $S_i^{k+1}$ is the set of $s_i \in S_i^k$ that no mixed strategy
$\sigma_i \in \Delta(S_i^k)$ strictly dominates on $S_{-i}^k$, i.e.\ with
$u_i(\sigma_i, s_{-i}) > u_i(s_i, s_{-i})$ for every $s_{-i} \in S_{-i}^k$. The survivors are
$\bigcap_k S_i^k$. Dominating strategies may be mixed (Tadelis \S6.3, pp.~114--117; the slides' IESDS
of Lecture~1).

\solution
Let $\sigma^* = (\sigma_1^*, \dots, \sigma_n^*)$ be a mixed-strategy Nash equilibrium and
$\operatorname{supp}\sigma_i^* = \{s_i : \sigma_i^*(s_i) > 0\}$.

\textbf{Claim.} $\operatorname{supp}\sigma_i^* \subseteq S_i^k$ for every player $i$ and every $k \ge 0$.

\textbf{Base, $k = 0$.} $S_i^0 = S_i$.

\textbf{Step, $k \to k+1$.} Assume $\operatorname{supp}\sigma_j^* \subseteq S_j^k$ for all $j$. Suppose
some $s_i \in \operatorname{supp}\sigma_i^*$ is not in $S_i^{k+1}$. Since $s_i \in S_i^k$ by assumption,
some $\sigma_i' \in \Delta(S_i^k)$ strictly dominates it on $S_{-i}^k$:
\[
  u_i(\sigma_i', s_{-i}) > u_i(s_i, s_{-i}) \quad\text{for every } s_{-i} \in S_{-i}^k .
\]
By the induction hypothesis the mixed profile $\sigma_{-i}^*$ puts probability one on $S_{-i}^k$.
Averaging the strict inequality with its weights $\sigma_{-i}^*(s_{-i})$, which are non-negative and
sum to one,
\begin{equation}
  u_i(\sigma_i', \sigma_{-i}^*) \;=\; \sum_{s_{-i} \in S_{-i}^k} \sigma_{-i}^*(s_{-i})\,u_i(\sigma_i', s_{-i})
  \;>\; \sum_{s_{-i} \in S_{-i}^k} \sigma_{-i}^*(s_{-i})\,u_i(s_i, s_{-i}) \;=\; u_i(s_i, \sigma_{-i}^*).
  \label{eq:gibbons-1-14-avg}
\end{equation}
Now use equilibrium. By the principle of indifference, a strategy in the support earns the equilibrium
payoff: $u_i(s_i, \sigma_{-i}^*) = u_i(\sigma^*)$. And since $\sigma^*$ is a Nash equilibrium,
$u_i(\sigma^*) \ge u_i(\sigma_i', \sigma_{-i}^*)$ for every $\sigma_i'$. Together,
\[
  u_i(s_i, \sigma_{-i}^*) = u_i(\sigma^*) \;\ge\; u_i(\sigma_i', \sigma_{-i}^*),
\]
which contradicts \eqref{eq:gibbons-1-14-avg}. Hence every $s_i \in \operatorname{supp}\sigma_i^*$ is in
$S_i^{k+1}$, which completes the induction. Since the claim holds for every $k$, the support lies in
$\bigcap_k S_i^k$. \hfill $\blacksquare$

\textbf{Where each assumption is used.} (i) Independence across players makes $\sigma_{-i}^*$ a
probability distribution on $S_{-i}$, so that the expectation \eqref{eq:gibbons-1-14-avg} is a weighted
average. (ii) The induction hypothesis puts that distribution inside $S_{-i}^k$, where the dominance
inequality holds; without it the inequality need not hold at the profiles that $\sigma_{-i}^*$ uses.
(iii) Indifference over the support is what turns ``$s_i$ is a best reply'' into
$u_i(s_i, \sigma_{-i}^*) = u_i(\sigma^*)$ for each $s_i$ in the support, not just on average.

\textbf{Why only the support.} Strategies outside the support may be eliminated or not; the
statement is about those played with positive probability. A strategy outside the support can be
strictly dominated, and the proof says nothing about it.

\textbf{Example.} In matching pennies no strategy is ever dominated, so IESDS removes nothing and
the support $\{H, T\}$ of the mixed equilibrium $(\frac12,\frac12)$ survives. In the game with payoffs
$u(a, L) = 3$, $u(a, R) = 0$, $u(b, L) = 0$, $u(b, R) = 3$, $u(c, \cdot) = 1$ for the row player, the
strategy $c$ is not strictly dominated by a pure strategy but is by the mix $\frac12 a + \frac12 b$
(payoff $\frac32 > 1$ against each column), so $c$ cannot be in the support of any equilibrium.

\interp Mixed equilibrium puts weight only on strategies that are rationalizable by iterated
dominance. So IESDS is a necessary filter for Nash equilibria in mixed as well as in pure strategies:
if IESDS leaves a single profile it is the unique Nash equilibrium (the converse direction of
Lecture~2, slide~17), and if a strategy is eliminated it never appears in the support of any equilibrium.

\verify \texttt{check\_L03.py} draws 300 random $3\times3$ bimatrix games with integer payoffs,
enumerates every equilibrium by support enumeration with exact fractions (checking each candidate
against all pure deviations), computes the IESDS survivors with strict dominance by mixed
strategies (linear programming) and asserts that every support lies inside the survivors; it also
checks the example above. Output line: \texttt{gibbons-1-14: 521 NE (119 with a mixed player) in 300
random games, every support inside the IESDS survivors}. Gibbons prints no answer.

\mistakes
\begin{itemize}
  \item Taking the dominating strategy $\sigma_i'$ to be pure: then the argument misses strategies
  dominated only by a mix.
  \item Forgetting the induction: dominance in the \emph{original} game is not the same as dominance
  in the reduced game $S^k$, and $\sigma_{-i}^*$ is known to live in $S_{-i}^k$ only because the
  opponents' supports have already been shown to survive $k$ rounds.
  \item Using only $u_i(\sigma^*) \ge u_i(s_i', \sigma_{-i}^*)$ for pure $s_i'$: the contradiction needs
  the equality $u_i(s_i, \sigma_{-i}^*) = u_i(\sigma^*)$ for the particular strategy $s_i$ in the support.
\end{itemize}

\answer{Induction on $k$: if all supports lie in $S^k$, a support strategy $s_i$ eliminated at round
$k+1$ would be strictly dominated by some $\sigma_i'$ on $S_{-i}^k$, hence against $\sigma_{-i}^*$:
$u_i(\sigma_i', \sigma_{-i}^*) > u_i(s_i, \sigma_{-i}^*) = u_i(\sigma^*)$, contradicting that $\sigma^*$ is a
Nash equilibrium. So all support strategies survive every round.}
\end{exsol}
